\subsection{CLAP}

In addition to the other models, we also try to use another audio feature extractor. The feature extractor we wish to try and use is one based from the CLAP model. CLAP, Contrastive Language-Audio Pretraining, is a model architecture which emerged back in November 2022\footnote{It should be noted that two major CLAP models exists (both have the exact same acronym) which were developed independently. These are from Microsoft and from LAION. To reduce confusion we will for the remainder of the thesis only refer to the CLAP model developed by LAION.\cite{elizalde_clap_2022}\cite{wu_large-scale_2024}}, where the audio domain was still dealing with challenges in supervised envionment due to the lack of large scale labeled data.

The contrastive part of the name refers to the learning objective for the model. In CLAP this refers to the Language-Audio, but can be for many groups of modalities, such as text-image in CLIP. The contrastive loss function utilized is simpliy a loss based on a sum of logged dot product between the models audio and text embeddings.

$$
    L = \frac{1}{2N} \sum_{i=1}^{N} \left( \log \frac{\exp(E_i^a\cdot E_i^t / \tau)}{\sum_{j=1}^{N} \exp(E_J^a \cdot E_j^t / \tau)} + \log \frac{\exp(E_i^t\cdot E_i^a / \tau)}{\sum_{j=1}^{N} \exp(E_J^t \cdot E_j^a / \tau)} \right)
$$

Contrastive learning is a paradigm used in machine learning to help models learn representations for data which is unlabeled. By feeding pairs of data into the model, it will help learn a shared latent space for audio and text. Perhaps these representations it has learned in such settings can be utilized in our pipeline.

The model was trained on audio files which were mono-channelled and at a sample rate of 48kHz.\cite{wu_large-scale_2024}
